\documentclass[11pt]{amsart}
\usepackage[T1]{fontenc}
\usepackage{lmodern}
\usepackage{amsmath,amssymb,amsthm}
\usepackage[margin=1.1in]{geometry}
\usepackage{booktabs}
\usepackage{enumitem}
\usepackage{graphicx}
\usepackage{hyperref}

\newtheorem{theorem}{Theorem}[section]
\newtheorem{lemma}[theorem]{Lemma}
\newtheorem{proposition}[theorem]{Proposition}
\newtheorem{corollary}[theorem]{Corollary}
\newtheorem{conjecture}[theorem]{Conjecture}
\theoremstyle{definition}
\newtheorem{definition}[theorem]{Definition}
\newtheorem{remark}[theorem]{Remark}
\newtheorem{example}[theorem]{Example}

\newcommand{\cD}{\mathcal D}
\newcommand{\cC}{\mathcal C}
\newcommand{\lean}[1]{\quad{\footnotesize[Lean: \texttt{#1}]}}
\newcommand{\TRN}{\mathrm{TRN}}
\newcommand{\NA}{\mathrm{NA}}
\emergencystretch=3em

\title[Conjecture N]{Conjecture N: the witness closes the hypothesis side.\\
A reduction to the tight transfer and an all-odd residual,\\
a protected Fan calculus, and the closure obstruction}
\author{Gallai project (cloud session \texttt{claude/conj-n})}
\date{27 November 2026 (working draft)}

\begin{document}
\begin{abstract}
Let $G$ have at most one vertex $g$ of positive even degree, let $x,y$ be odd and non-adjacent,
and let $y$ have a neighbour $w\ne g$ with $w\notin N(x)$. Conjecture~N of~\cite{DL} asks for a
Lov\'asz decomposition whose path $Q(y)$ ending at $y$ avoids $x$. With Proposition~9.2
of~\cite{DL} it gives DL(1). Unlike the avoidance conjecture $L_1^*$, it allows $x\sim g$.
We ran the counterexample-guided search for an inductive class of avoidance instances that
contains every instance of N and is closed under the proved transfers. We did not find a closed
class. What we prove and find is the following.
\begin{itemize}
\item \emph{Proved and formalised} (Lean, no \texttt{sorry}, standard axioms): \textbf{N follows from
two statements}, $\TRN\wedge\NA\Rightarrow$ N (Theorem~\ref{thm:red}).
  \begin{itemize}
  \item $\TRN$ is the tight transfer conjecture TR of~\cite{L1I} for one forbidden vertex, with the
  tail also allowed to be adjacent to $g$.
  \item $\NA$ is a statement about graphs with \emph{no} even vertex: $Q(y)$ can avoid a
  \emph{neighbour} $x$ when $y$ has another neighbour not adjacent to $x$.
  \item $\NA$ is exactly the case $N(g)=\{x,y\}$ of N.
  \end{itemize}
  The induction never leaves the class of N. A non-tight target always gives a child that is again
  an instance of N, because the witness survives. So the hypothesis side of Lov\'asz's Case~I,
  which is the obstruction for $L_1^*$ (the conjecture ST of~\cite{L1I}), is automatic here. The
  proof uses no $L_1^*$. For the instances of the blocked case of $L_1^*$ with $|A|=1$ that N
  covers, it replaces ST by the all-odd statement NA.
\item \emph{Proved}: three new transfers.
  \begin{itemize}
  \item A leaf-path move.
  \item A \emph{protected} form of Fan's lemma, which keeps one path of the old decomposition
  unchanged.
  \item From it, moves for instances with a forbidden set $B$ of far ends of $Q(c)$. Among them is
  a \emph{target-in-$A$} move, Case~I across a target that lies in the forbidden set.
  \end{itemize}
  The protected Fan lemma, in its basic form, and the target-in-$A$ move are also formalised.
\item \emph{Computed}, all exhaustive in the stated range:
  \begin{itemize}
  \item With the calculus of~\cite[Thm.~3.1]{L1I}, $115$ of the $263{,}264$ instances of N of
  order $9$ have no certificate. With the new moves, $7$ remain; all have $g$ of degree two on a
  triangle.
  \item $\TRN$: no violation in $490{,}923$ tests of order $9$ and $831{,}785$ on two order-$11$
  slices.
  \item NA: no failure through order $10$ ($1{,}305{,}621$ instances at order $10$).
  \item Under $\TRN$, the class N $\cup$ NA has \emph{no} stuck instance at orders $5$, $7$, $9$.
  \end{itemize}
\item \emph{The obstruction.}
  \begin{itemize}
  \item Under the unconditional calculus, N is stuck exactly at \emph{all-tight} configurations,
  where every neighbour of $g$ other than $x$ is $y$ or adjacent to $y$ ($37{,}652$ of the
  $263{,}264$ at order $9$).
  \item Every enlargement we tried either stays stuck or contains false instances: set forms,
  witness forms, far-end forms, subcubic graphs.
  \item The persistent shapes are (i) the all-tight configurations, which need a transfer through
  a tight target; (ii) NA, whose leaf-graph children carry a forbidden pair $\{x,g\}$ with
  $x\in N(g)$; and (iii) an even vertex of degree two on a triangle, a subdivided double edge.
  \end{itemize}
\end{itemize}
Conjecture~N itself is not proved here.
\end{abstract}
\maketitle

\section{Introduction}\label{sec:intro}

All graphs are finite and simple. A vertex is \emph{even} if its degree is positive and even, and
\emph{odd} if its degree is odd. A \emph{Lov\'asz decomposition} of a graph with at most one even
vertex partitions its edges into paths so that every vertex $v$ is an end of exactly $d(v)\bmod2$
of them~\cite{L68}. For an odd vertex $c$, $Q(c)=Q_c$ is the path ending at $c$, and its other end is
the \emph{far end} of $Q_c$. Throughout, $g$ is the even vertex when there is one.

\begin{conjecture}[N; {\cite[Conjecture~9.1]{DL}}]\label{conj:n}
Let $G$ have at most one even vertex $g$, let $x,y$ be odd and non-adjacent, and suppose that the
\emph{witness set}
\[
W(G,y,x)=N(y)\setminus\bigl(N(x)\cup\{g\}\bigr)
\]
is non-empty. Then some Lov\'asz decomposition of $G$ has $x\notin V(Q(y))$.
\lean{ConjNOn}
\end{conjecture}

\subsection*{Why it matters}
\begin{itemize}
\item By~\cite[Prop.~9.2]{DL}, N reduces DL(1) to \emph{near twins}, pairs with
$N(a)\setminus\{g\}=N(b)\setminus\{g\}$.
\item DL(1), and its multi-pair form DL$(k)$, gives the two-end test (T1) at every root whose data
have $M$ connected (\texttt{test1At\_of\_disjL1}).
\item N contains most of the blocked case of $L_1^*$ with $|A|=1$~\cite[Remark~9.3]{DL}, and also
the instances with $x\sim g$, which $L_1^*$ excludes.
\item No counterexample is known: there are $264{,}437$ instances of order at most $9$, all good.
\end{itemize}

\subsection*{The method}
We look for a class $\cC$ of avoidance instances $(G,c,A)$, defined by a local predicate, such that
\begin{enumerate}
\item $\cC$ contains every instance of N;
\item $\cC$ consists of good instances;
\item every instance of $\cC$ is a base case, or has a proved move to a strictly smaller instance of
$\cC$.
\end{enumerate}
The loop is: propose $\cC$; compute its \emph{closure failures} (instances of $\cC$ with no move
into $\cC$) exhaustively in small orders; refine. A class with no failures would be an inductive
hypothesis for N, and a hand proof of its closure would prove N. The tool is
\texttt{closure2.cpp}, the oracle is the exact solver of~\cite{L1I}, and Section~\ref{sec:comp}
lists the runs.

\subsection*{Results}
\begin{enumerate}
\item Section~\ref{sec:moves}, \emph{three new transfers}, proved on paper:
  \begin{itemize}
  \item the leaf-path move (Lemma~\ref{lem:leafpath});
  \item Fan's lemma keeping a protected path (Lemma~\ref{lem:protected});
  \item the far-end calculus (Theorem~\ref{thm:far}), with the target-in-$A$ move~(x).
  \end{itemize}
\item Section~\ref{sec:red}, \emph{the reduction}, proved and formalised:
  \begin{itemize}
  \item $\TRN\wedge\NA\Rightarrow$ N (Theorem~\ref{thm:red});
  \item the degree-two case of N is NA (Lemma~\ref{lem:deg2});
  \item the unconditional part (Proposition~\ref{prop:uncond}): a smallest counterexample to N is
  all-tight.
  \end{itemize}
\item Section~\ref{sec:comp}, \emph{computations}: certificates, data for $\TRN$, NA and the
strengthening $N^+$, and closure tables.
\item Section~\ref{sec:obs}, \emph{the obstruction}, with the smallest example of each failing
refinement.
\item Section~\ref{sec:lean}, \emph{formalisation}.
\end{enumerate}

\subsection*{Status}
\begin{itemize}
\item \emph{Proved and formalised}:
  \begin{itemize}
  \item Theorem~\ref{thm:red}, Lemma~\ref{lem:deg2}, Lemma~\ref{lem:split},
  Corollary~\ref{cor:TR};
  \item Lemma~\ref{lem:protected}, in the case where no end of $P$ lies in $N_K(u)$;
  \item move~(x) of Theorem~\ref{thm:far}.
  \end{itemize}
\item \emph{Proved on paper}:
  \begin{itemize}
  \item Lemma~\ref{lem:leafpath};
  \item Lemma~\ref{lem:protected} with ends of degree one in $N_K(u)$;
  \item the moves (b$_B$), (C), (c$_B$) of Theorem~\ref{thm:far};
  \item Proposition~\ref{prop:uncond}, whose content is also the first step of the formal proof.
  \end{itemize}
\item \emph{Verified}: by exact computation, in the stated range.
\item \emph{Conjectured}: supported by the data quoted, not proved.
\end{itemize}
No proof uses a computation.

\section{Instances, moves and base cases}\label{sec:prelim}

For a path decomposition $\cD$, $b_\cD(v)$ is the number of paths of $\cD$ with an end at $v$, and
$Z_\cD(u)=\{w\in N(u):b_\cD(w)=0\}$ is the set of \emph{blanks} of $u$.

\begin{definition}
An \emph{instance} $(G,c,A)$ has $G$ with at most one even vertex, $c$ odd, and $c\notin A$. It is
\emph{good} if some Lov\'asz decomposition has $V(Q_c)\cap A=\emptyset$. An \emph{extended
instance} $(G,c,A,B)$ has in addition a set $B$ of odd vertices. It is good if some Lov\'asz
decomposition has $V(Q_c)\cap A=\emptyset$ and the far end of $Q_c$ outside $B$. A \emph{target}
is a vertex $u\in N(g)$; it is \emph{tight} if $u\in N(c)$.
\end{definition}

An instance of N is $(G,y,\{x\})$ under the hypotheses of Conjecture~\ref{conj:n}; we write it
$(G,y,x)$.

\begin{lemma}[Fan's lemma with end paths~{\cite{Fan05}, \cite[Lemma~2.2]{RD}, \cite[Thm.~5.1]{CLT}, \cite[Lemma~2.1]{DL}}]\label{lem:fan}
Let $t\ne u$ be non-adjacent vertices of $K$ and $\cD$ a decomposition of $K$ with
$b_\cD(t)>|Z_\cD(u)|$. Then $K+tu$ has a decomposition $\cD^+$ with
$b_{\cD^+}=b_\cD-\mathbf 1_t+\mathbf 1_u$. Every path $Q$ of $\cD^+$ comes from a path $Q'$ of $\cD$
with $V(Q)\subseteq V(Q')\cup\{u\}$, and every end $w\notin\{t,u\}\cup N_K(u)$ of $Q$ is an end of
$Q'$. If $b_\cD\le1$, the chain is \emph{forced}: $x_0=t$, and $x_{i+1}$ is the vertex before $u$
on the path ending at $x_i$, read from $x_i$. Only the paths ending at $x_0,\dots,x_k$ change. The
first $k$ are re-routed, which keeps their vertex sets and replaces the end $x_i$ by $x_{i+1}$; then
$u$ is appended to the last.
\end{lemma}

\begin{theorem}[the transfer calculus~{\cite[Thm.~3.1]{L1I}}]\label{thm:calc}
For an instance $(G,c,A)$, goodness of the child on the right implies goodness of $(G,c,A)$:
\begin{enumerate}[label=(\alph*)]
\item $c\in N(g)$, $g\notin A$: $(G-gc,\ g,\ A\cup\{c\})$;
\item $u\in N(g)\setminus(A\cup N[c])$: $(G-gu,\ c,\ A)$;
\item $u\in N(g)\cap N(c)\setminus A$: $(G-cu-gu,\ u,\ A\cup\{c\})$;
\item no even vertex, $d(c)\ge3$, $y\in N(c)\setminus A$: $(G-cy+y\ell,\ \ell,\ A\cup\{c\})$, a new leaf $\ell$;
\item $d(g)=2$, $N(g)=\{u,v\}$, $uv\notin E$, $G'=G-gu-gv+uv$: $(G',c,A)$ if $g\notin A$, and
$(G',c,A-g+w)$ for $w\in\{u,v\}\setminus\{c\}$ if $g\in A$;
\item $d(g)=2$, $N(g)=\{u,v\}$, $uv\in E$, $g\in A$: $(G,c,A\setminus\{g\})$;
\item $d(c)=3$, $y_1,y_2\in N(c)$ non-adjacent: $(G-cy_1-cy_2+y_1y_2,\ c,\ A\cup\{y_1\})$;
\item $g\notin A$, $R$ a $g$--$c$ path with interior outside $A$: $(G-E(R),\ g,\ A\cup V(R)\setminus\{g\})$;
\item $C$ the component of $c$ in $G-A$: $(G,c,\partial C)$, where $\partial C$ is the set of
vertices of $A$ with a neighbour in $C$.
\end{enumerate}
\end{theorem}

The base cases are as in~\cite{L1I,DL}:
\begin{itemize}
\item $A$ misses the component of $c$;
\item no even vertex, $d(c)=1$, and the neighbour of $c$ is not in $A$;
\item $A\cup\{c\}$ packed, by separation~\cite[Thm.~6.1]{RD};
\item no even vertex, $A$ packed and disjoint from $N[c]$ ($L_1$, \texttt{lovasz\_avoid\_set});
\item every vertex has at most one neighbour in $A\cup\{c\}$, and $A\cap N[c]=\emptyset$; this
is \texttt{csys\_good\_of\_sparse};
\item the leaf case $N(x)=\{g\}$ (\texttt{avoid\_leaf}), and the degree-two case
(\texttt{avoid\_deg\_two}) of~\cite{DL}.
\end{itemize}
For extended instances with $B\ne\emptyset$ only two base cases are used:
\begin{itemize}
\item $A\cup B$ misses the component of $c$;
\item no even vertex, $d(c)=1$, and the neighbour of $c$ is not in $A\cup B$. Here $Q_c=cv$: take
a Lov\'asz decomposition of $G-cv$ and add $cv$.
\end{itemize}

\section{Three new transfers}\label{sec:moves}

\begin{lemma}[leaf paths]\label{lem:leafpath}
Let $P$ be a path of $G$ whose two ends have degree one in $G$, and let $c$ not be an end of $P$.
If $(G-E(P),c,A,B)$ is good, so is $(G,c,A,B)$.
\end{lemma}

\begin{proof}
Interior vertices of $P$ lose two edges, so they keep their parity, and the ends of $P$ become
isolated. So $G-E(P)$ has at most one even vertex and $c$ is odd in it. Add $P$ to a good
decomposition of $G-E(P)$. The two ends of $P$ get one end each, as their degree one requires, and
nothing else changes. So the result is a Lov\'asz decomposition of $G$, with the same $Q_c$.
\end{proof}

\begin{lemma}[Fan's lemma with a protected path]\label{lem:protected}
Let $t\ne u$ be non-adjacent vertices of $K$, and let $\cD$ be a decomposition of $K$ with
$b_\cD(t)>|Z_\cD(u)|$. Let $P\in\cD$ be a path. Suppose that each end $e$ of $P$ satisfies
$e\notin\{t,u\}$, and either $e\notin N_K(u)$ or $d_K(e)=1$. Then $K+tu$ has a decomposition
$\cD^+$ with $P\in\cD^+$ and $b_{\cD^+}=b_\cD-\mathbf 1_t+\mathbf 1_u$.
\end{lemma}

\begin{proof}
Let $e_1,e_2$ be the ends of $P$, $K^\circ=K-E(P)$ and $\cD^\circ=\cD\setminus\{P\}$. Then
$b_{\cD^\circ}=b_\cD-\mathbf 1_{e_1}-\mathbf 1_{e_2}$.
\begin{itemize}
\item $t$ is not an end of $P$, so $b_{\cD^\circ}(t)=b_\cD(t)$.
\item Let $w\in Z_{\cD^\circ}(u)$, so $w\in N_{K^\circ}(u)$. Suppose $w$ were an end of $P$. Then
$w\in N_K(u)$, so $d_K(w)=1$. Its only edge lies on $P$, so $w$ is isolated in $K^\circ$, a
contradiction. Hence $b_\cD(w)=b_{\cD^\circ}(w)=0$ and $w\in Z_\cD(u)$.
\end{itemize}
So $b_{\cD^\circ}(t)>|Z_{\cD^\circ}(u)|$, and $t\not\sim u$ in $K^\circ$. Lemma~\ref{lem:fan}
gives a decomposition of $K^\circ+tu$ with $b=b_{\cD^\circ}-\mathbf 1_t+\mathbf 1_u$. Adding $P$,
whose edges are not in $K^\circ+tu$, gives $\cD^+$.
\end{proof}

Lemma~\ref{lem:fan} alone does not give this. Its tracking statement allows the path coming from
$P$ to gain $u$ or to change its ends; Lemma~\ref{lem:protected} keeps $P$ itself.

\begin{theorem}[the far-end calculus]\label{thm:far}
Let $(G,c,A,B)$ be an extended instance with even vertex $g$. In each case below, goodness of the
child implies goodness of $(G,c,A,B)$.
\begin{enumerate}[label=(\roman*)]
\item[(x)] \emph{Target in $A$.} $u\in N(g)\setminus N[c]$, \emph{with $u$ allowed in $A$ or $B$}:
child $(G-gu,\ c,\ A,\ B\cup N_G(u))$.
\item[(b$_B$)] $u\in N(g)\setminus(A\cup B\cup N[c])$ and $N(u)\cap B=\emptyset$: child
$(G-gu,\ c,\ A,\ B)$.
\item[(C)] $u\in N(g)\cap N(c)\setminus A$: child $(G-cu-gu,\ u,\ A\cup\{c\},\ B\cup N_G(u))$.
\item[(c$_B$)] $u\in N(g)\cap N(c)\setminus A$ and $N_G(u)\cap B=\emptyset$: child
$(G-cu-gu,\ u,\ A\cup\{c\},\ B)$.
\end{enumerate}
The moves (a), (d)--(h) and the leaf-path move carry $B$ unchanged. Move (i) carries $B$ as
$B\cap C$.
\end{theorem}

\begin{proof}
(x). Let $G'=G-gu$ and let $\cD'$ be good for the child, so $Q'_c$ avoids $A$ and its far end $q$
is not in $B\cup N_G(u)$. In $G'$:
\begin{itemize}
\item $g$ is odd and has one end;
\item the only possible even vertex is $u$, so every neighbour of $u$ is odd, has one end, and
$Z(u)=\emptyset$.
\end{itemize}
The ends of $Q'_c$ satisfy the hypothesis of Lemma~\ref{lem:protected} with $t=g$:
\begin{itemize}
\item $c\notin\{g,u\}\cup N_{G'}(u)$, since $u\notin N[c]$;
\item $q\notin N_G(u)=N_{G'}(u)\cup\{g\}$, and $q\ne u$, since $u$ has no end in $G'$.
\end{itemize}
Lemma~\ref{lem:protected} gives a decomposition of $G$ that contains $Q'_c$ and moves the end of $g$
to $u$; it is a Lov\'asz decomposition. So $Q_c=Q'_c$.

(b$_B$). Apply move (b) through Lemma~\ref{lem:fan}. The new $Q_c$ comes from a path with an end at
$c$, because $c\notin\{g,u\}\cup N_{G'}(u)$; that path is $Q'_c$. So $V(Q_c)\subseteq
V(Q'_c)\cup\{u\}$, which avoids $A$. Let $f$ be the far end of $Q_c$. If $f\notin\{g,u\}\cup
N_{G'}(u)$, it is an end of $Q'_c$ other than $c$, so $f=q\notin B$. Otherwise $f\in N_G[u]$, which
misses $B$.

(C). Let $K=G-cu-gu$. Its only possible even vertex is $c$. Let $\cD_K$ be good for the child:
$R_0=Q^K_u$ avoids $A\cup\{c\}$, and its far end $p$ is not in $B\cup N_G(u)$.
\begin{itemize}
\item \emph{Leaf graph.} Add a new leaf $\ell$ at $u$, standing for $c$, and extend $R_0$ to
$R=\ell uR_0$. In $H=K+u\ell$ the vertex $u$ is even with no end. The neighbours of $u$ other than
$\ell$ are odd with one end, and $\ell$ has an end, so $Z(u)=\emptyset<1=b(g)$.
\item \emph{Fan step.} The ends of $R$ are $\ell$ and $p$. Here $\ell\in N_H(u)$ has degree one,
and $p\notin\{g,u\}\cup N_H(u)$. Lemma~\ref{lem:protected} with $t=g$ gives a decomposition of
$H+gu$ containing $R$.
\item \emph{Merge.} $H+gu$ is $G-cu$ with the leaf $\ell$ at $u$ standing for $c$. Only $R$
contains $\ell$, and $R$ avoids $c$. So $\ell$ merges into $c$~\cite[Lemma~12]{MS}, and $c$
receives exactly the end of $\ell$.
\end{itemize}
The result is a Lov\'asz decomposition of $G$ with $Q_c=cuR_0$. This path avoids $A$, and its far
end is $p\notin B$.

(c$_B$). Make the same construction, with Lemma~\ref{lem:fan} in place of
Lemma~\ref{lem:protected}. In $H$ every vertex has at most one end, so the chain is forced.
\begin{itemize}
\item $\ell$ is never a chain vertex: it precedes $u$ only on $R$ read from $\ell$, and $x_0=g\ne\ell$.
\item $R$ contains $u$, so it is never appended to. It changes only if $p=x_i$ for some $i$. Then it
is re-routed once: it keeps its vertex set, and its end $p$ becomes $x_{i+1}$. This vertex is the
predecessor of $u$ on $R$ read from $p$, so it lies in $N_H(u)\setminus\{\ell\}\subseteq N_G(u)$.
\end{itemize}
So after the merge, $Q_c$ has vertex set $V(R_0)\cup\{c\}$, and its far end lies in $\{p\}\cup
N_G(u)$, which misses $B$.

For the remaining moves, the proofs of~\cite[Thm.~3.1]{L1I} keep the far end. In (a), (h) the new
$Q_c$ ends at the far end of $Q'_g$. In (d) the leaf is merged and the far end is unchanged. In
(e), (f), (g) and the leaf path, $Q_c$ is unchanged or subdivided. In (i) a path from $c$ that
avoids $\partial C$ stays inside $C$.
\end{proof}

\begin{remark}
Move~(x) is the mechanism of~\cite[Prop.~9.5]{DL}. That conjecture (SF) is the statement that the
(x)-child at the target $a$ is good. Move (x) is unconditional, but its child carries the
far-end set $N_G(u)$. Moves (x) and (C) are the only ones that create far-end constraints.
\end{remark}

\section{The reduction}\label{sec:red}

\begin{conjecture}[$\TRN$]\label{conj:trn}
Let $g$ be the only even vertex of $G$, let $x,y$ be odd and non-adjacent, and let
$u\in N(g)\cap N(y)$ with $u\ne x$. If $(G-gu,y,\{x\})$ is good, then $(G,y,\{x\})$ is good.
\lean{TightTransferNOn}
\end{conjecture}

For $y\notin N(g)$, this is the tight transfer conjecture TR of~\cite[Conj.~4.1]{L1I} with $A=\{x\}$.
For $y\in N(g)$ it is the same statement for a tail adjacent to $g$. We write TR$^*$ for TR with
the hypothesis $c\notin N(g)$ dropped.

\begin{conjecture}[NA, neighbour avoidance]\label{conj:na}
Let $G$ have no even vertex, let $x\sim y$, and let $y$ have a neighbour $w\ne x$ with
$w\notin N(x)$. Then some Lov\'asz decomposition has $x\notin V(Q(y))$.
\lean{NeighbourAvoidOn}
\end{conjecture}

\begin{lemma}[splitting off a vertex of degree two]\label{lem:split}
Let $N(g)=\{a,b\}$ with $a\ne b$ and $ab\notin E$, and let $H=G-ga-gb+ab$. If $c\ne g$, $g\notin
A$, and $H$ has a Lov\'asz decomposition whose path ending at $c$ avoids $A$, then so has $G$.
\lean{avoid\_of\_splitOff}
\end{lemma}

\begin{proof}
Subdivide $ab$ on its path by $g$. The ends are unchanged, $g$ is interior of degree two, and
$Q_c$ gains at most $g$.
\end{proof}

\begin{lemma}[the degree-two case is NA]\label{lem:deg2}
Let $(G,y,x)$ be an instance of N with $N(g)=\{x,y\}$. Then $H=G-gx-gy+xy$ has no even vertex,
$x\sim y$ in $H$, and $W(G,y,x)$ consists of neighbours of $y$ in $H$ not adjacent to $x$. So NA
for $(H,y,x)$ implies N for $(G,y,x)$. Conversely, subdividing $xy$ turns every instance of NA into
an instance of N with $N(g)=\{x,y\}$. NA holds when $x,y$ have no common neighbour; for N with
$N(g)=\{x,y\}$ this is the case $N(x)\cap N(y)=\{g\}$, covered by~\cite[Thm.~6.1]{DL}.
\lean{conjNDegTwo\_of\_neighbourAvoid}
\end{lemma}

\begin{proof}
In $H$, $g$ is isolated and every other vertex keeps its degree. The new edge is $xy$. A witness
$w\ne g$ is not adjacent to $x$ in $H$, since $w\ne y$. Apply NA in $H$, then
Lemma~\ref{lem:split}; $Q(y)$ avoids $x$, so it does not use $xy$. For the converse, the
subdivided graph has the single even vertex $g$ of degree two. For the last claim,
take~\cite[Thm.~6.1]{DL} with $a=y$, $b=x$; $d(y)\ge3$, since $y\sim g$ and $y\sim w$.
\end{proof}

\begin{theorem}\label{thm:red}
$\TRN$ and the degree-two case of N together imply Conjecture~N. Hence $\TRN\wedge\NA$ implies N.
\lean{conjN\_of\_tight\_degTwo, conjN\_of\_tight\_neighbourAvoid}
\end{theorem}

\begin{proof}
Induction on $\sum_vd(v)$. Let $(G,y,x)$ be an instance of N, and fix a witness $w_0\in W(G,y,x)$.
\begin{enumerate}[label=(\arabic*)]
\item \emph{No even vertex.} $L_1$ (\texttt{lovasz\_avoid\_nonadj}) gives $Q(y)\not\ni x$. From now
on $g$ exists. Note $x,y\ne g$, $w_0\ne x,y,g$.
\item \emph{A non-tight target.} Suppose $u\in N(g)\setminus\{x,y\}$ and $u\not\sim y$. In $G-gu$
the only possible even vertex is $u$; the degrees of $x,y$ are unchanged; $x\not\sim y$. Also
$w_0\ne u$, as $w_0\in N(y)$, so $w_0$ is a witness of the child. By induction the child is good,
and Case~I with a target outside $N[y]\cup\{x\}$ transfers this (\texttt{csys\_caseI},
\texttt{avoid\_free\_target}). This holds whether or not $y\sim g$.
\item \emph{A tight target with a witness.} Suppose $u\in N(g)\cap N(y)$, $u\ne x$, and some $w\in
N(y)\setminus N(x)$ has $w\ne u$. In $G-gu$ the even vertex is $u$, and $w\ne u$ is a witness; it
may be the old $g$, now odd. By induction $(G-gu,y,x)$ is good, and $\TRN$ gives $(G,y,x)$.
\item \emph{The rest.} Every $u\in T:=N(g)\setminus\{x,y\}$ is adjacent to $y$, by (2). By (3), for
every $u\in T$ every $w\in N(y)\setminus N(x)$ equals $u$.
  \begin{itemize}
  \item Applied to $w_0$: $T\subseteq\{w_0\}$.
  \item Applied to $w=g$ when $w_0\in N(g)$: if $y\sim g$ then $x\sim g$, since $g\ne w_0$.
  \end{itemize}
  So $N(g)\subseteq\{x,y,w_0\}$, and $d(g)$ is positive and even.
  \begin{itemize}
  \item If $w_0\notin N(g)$, then $N(g)=\{x,y\}$, the degree-two case.
  \item If $w_0\in N(g)$, then $N(g)\ne\{x,y,w_0\}$, which has odd size. Also $N(g)\ne\{w_0\}$.
  The set $\{w_0,y\}$ would give $y\sim g$ but $x\not\sim g$, which is excluded. So
  $N(g)=\{x,w_0\}$ with $x\not\sim w_0$. Split $g$ off: $H=G-gx-gw_0+xw_0$ has no even vertex and
  $x\not\sim y$ in it, so $L_1$ gives a good decomposition of $H$, and Lemma~\ref{lem:split}
  finishes.\qedhere
  \end{itemize}
\end{enumerate}
\end{proof}

\begin{corollary}\label{cor:TR}
TR of~\cite{L1I} ($y\notin N(g)$), its counterpart for tails adjacent to $g$, and NA imply~N.
\lean{tightTransferNOn\_of, conjN\_of\_TR\_neighbourAvoid}
\end{corollary}

\begin{proposition}[the unconditional part]\label{prop:uncond}
\begin{itemize}
\item A smallest counterexample to~N is \emph{all-tight}: every neighbour of $g$ other than $x$ is
$y$ or a neighbour of $y$.
\item Without $\TRN$, step~(2) of the proof, with steps (1) and (4) in the cases that do not use
$\TRN$, proves N for every instance from which repeated non-tight targets lead to an instance with
no even vertex, or to $N(g)=\{x,u\}$ with $x\not\sim u\in W$.
\end{itemize}
\end{proposition}

\begin{remark}[N and the blocked case of $L_1^*$]
\begin{itemize}
\item \emph{The setting.} Let $(G,c,\{x\})$ be blocked in the sense of~\cite{L1S,L1I}: $c\notin
N(g)$, $x\notin N[c]\cup N(g)$, and $N(g)\subseteq N(x)$. Lov\'asz's Case~I at a non-tight target
$u$ produces the child $(G-gu,c,\{x\})$. That child has $x\sim u$, the new even vertex, so it is
not admissible for $L_1^*$. This is the \emph{hypothesis side} of~\cite{L1I}, the conjecture ST.
\item \emph{What N changes.} If $c$ has a neighbour outside $N(x)$, the child is again an instance
of N, and step~(2) applies. So for such instances, Theorem~\ref{thm:red} replaces ST by NA, a
statement about graphs without an even vertex. The transfer side TR is still needed.
\item \emph{The other gain.} The induction runs inside N, where $x\sim g$ is allowed, and never
calls $L_1^*$.
\end{itemize}
\end{remark}

\section{Computations}\label{sec:comp}

\emph{This section is computational; no proof above depends on it.}
\begin{itemize}
\item Scripts: \texttt{manuscripts/conj-n/scripts}. The README there gives the commands.
\item Graphs: all connected graphs with at most one even vertex (\texttt{nauty-geng -c}, filtered),
and the lists \texttt{one\_even\_\{5,7,9\}.g6} of~\cite{STL}.
\end{itemize}

\subsection*{Tools and validation}
\begin{itemize}
\item \emph{Oracle.} Goodness of $(G,c,A)$ is decided by the exact oracle of~\cite{L1I}
(\texttt{lind.h}). For extended instances, \texttt{orcb.h} adds the far-end restriction to the
first path.
\item \emph{Validation of the far-end oracle.} It agrees with full enumeration of Lov\'asz
decompositions (\texttt{enum.c} of~\cite{DL}) on all $399{,}762$ queries $(c,A,B)$ of the graphs of
orders $4$--$6$ with at most one even vertex and the $130$ graphs of order~$7$
(\texttt{xcheckb.py}). With $B=\emptyset$ it agrees with \texttt{lind.h} on every child tested.
\item \emph{Moves.} \texttt{ext.h} implements Theorem~\ref{thm:calc}, Lemma~\ref{lem:leafpath}
(paths of length two), Theorem~\ref{thm:far}, the base cases of Section~\ref{sec:prelim}, and
optionally the move (t) of TR or TR$^*$.
\item \emph{Closure.} \texttt{closure2.cpp} enumerates every instance of a class, every $A$, and
every $B$ when the class allows it, and reports \texttt{STUCK} and \texttt{FALSE} instances.
\end{itemize}

\subsection{Certificates}
An instance is \emph{certified} if a sequence of moves leads to a base case.
\begin{itemize}
\item \emph{Old calculus.} With Theorem~\ref{thm:calc} and the bases of~\cite{L1I,DL}
(\texttt{ncert.cpp}), all $12+1{,}161$ instances of N of orders $5,7$ are certified, but $115$ of the
$263{,}264$ of order $9$ are not. All $115$ are good. So no class closes under the old calculus.
\item \emph{New calculus.} With the moves of Section~\ref{sec:moves} (\texttt{xcert.cpp}), $108$ of
the $115$ are certified.
\item \emph{The $7$ left} (\texttt{data/xcert9\_uncertified.txt}) all have $d(g)=2$ with
$N(g)=\{u,x\}$, $u\sim x$, $u\in N(y)$, $y\notin N(g)$ and $d(y)=5$. An example is
\texttt{HsOceO\}}, $g=5$, $y=0$, $x=8$.
\end{itemize}

\subsection{Data for the conjectures}
All runs are exhaustive over the stated graphs. The outputs are in
\texttt{data/conjecture\_tests.txt}.
\begin{center}\small
\resizebox{\textwidth}{!}{\begin{tabular}{lll}
\toprule
statement & range & result\\
\midrule
$\TRN$, $y\notin N(g)$ / $y\in N(g)$ & order $9$ & $236{,}891$ / $254{,}032$ tests, $0$ violations\\
$\TRN$, both & order $11$, slices $17$, $517$ of $2{,}000$ & $831{,}785$ tests, $0$ violations\\
TR$^*$ with $c\in N(g)$, $|A|\le2$ & order $9$ (all $A$ at order $7$) & $565{,}744$ ($2{,}036$) tests, $0$ violations\\
NA & orders $4,6,8,10$ & $3$, $129$, $4{,}970$, $1{,}305{,}621$ instances, all good\\
$N^+$: $Q(y)\not\ni x$, far end $\notin N(x)$, for $x\sim g$ & order $9$ & $134{,}911$ of $134{,}911$\\
$N^+$ for $x\not\sim g$ & order $9$ & fails in $417$ of $128{,}353$\\
(x)-child of the core & order $9$ & $18{,}315$ of $18{,}315$ good\\
\bottomrule
\end{tabular}}
\end{center}
\begin{itemize}
\item The \emph{core} consists of the instances of N with $x\sim g$ and $N(g)\setminus\{x,y\}\subseteq
N(x)$. Its (x)-child is $(G-gx,y,\{x\},N_G(x))$.
\item Over all $134{,}911$ instances with $x\sim g$, the (x)-child is good in $134{,}869$; the $42$
failures are outside the core.
\item The far-end form of ``$Q(c)$ avoids $g$'', $(G,c,\{g\},N(g))$, fails in $85$ of $46{,}638$
cases at order~$9$, although $(G,c,\{g\})$ is good in all of them.
\end{itemize}

\subsection{Closure tables}
The instances are those of the class at the given order, with every $A$, and every $B$ for the
far-end classes. ``Stuck'' counts instances with no move into the class or a base case; ``false''
counts instances that are not good. The classes are:
\begin{itemize}
\item N$_1$: Conjecture~N, $|A|=1$.
\item N$_2$: the set form of~\cite{DL}: $A\cap N[c]=\emptyset$, and some $w\in N(c)\setminus\{g\}$
has $N(w)\cap A\subseteq\{g\}$.
\item V: as N$_2$, but $w=g$ is allowed.
\item L$_1$: $L_1^*$-admissible, $A\cap(N[c]\cup N(g))=\emptyset$.
\item U $=$ N$_2\cup$ L$_1$.
\item C$_1$ $=$ N$_1\cup$ NA $\cup\,($N$_2$ with $A\subseteq N[g])$.
\item Far-end classes:
  \begin{itemize}
  \item N$_2^B$: N$_2$ with any $B$, the witness outside $B$;
  \item N$_2^{N(A)}$: the same with $B\subseteq N(A)$;
  \item XC: N$_2$ together with the (x)-children of N$_2$ instances.
  \end{itemize}
\item W: a witness $w\notin A\cup\{g\}$ with $N(w)\cap A\subseteq\{g\}$, while $A$ may meet $N(c)$.
\end{itemize}
The moves are:
\begin{itemize}
\item U: the unconditional calculus of Section~\ref{sec:moves};
\item T: with TR ($c\notin N(g)$);
\item T$^*$: with TR$^*$.
\end{itemize}
Source: \texttt{data/tables.txt}.

\begin{center}\small
\resizebox{\textwidth}{!}{\begin{tabular}{llrrrr}
\toprule
class & moves & order 5: stuck/size & order 7: stuck/size & order 9: stuck/size & false\\
\midrule
N$_1$ & U & $4/12$ & $214/1{,}161$ & $37{,}652/263{,}264$ & $0$\\
N$_1$ & T & $4/12$ & $154/1{,}161$ & & \\
N$_1$ & T$^*$ & $0/12$ & $14/1{,}161$ & $1{,}521/263{,}264$ & \\
N$_1\cup$NA & T$^*$ & $0/12$ & $0/1{,}161$ & $0/263{,}264$ & \\
\quad (NA part, orders 6, 8) & T$^*$ & \multicolumn{3}{l}{order 6: $22/215$;\ order 8: $2{,}522/9{,}558$} & \\
C$_1$ & T$^*$ & $0/44$ & $164/2{,}842$ & \multicolumn{1}{l}{(orders 6, 8: $0$)} & $0$\\
N$_2$ & U & $2/46$ & $249/3{,}638$ & $112{,}226/1{,}066{,}750$ & $0$\\
N$_2$ & T / T$^*$ & $2$ / $0$ & $173$ / $72$ & & \\
N$_2$, subcubic & U & $2/36$ & $8/812$ & $182/22{,}592$ & $0$\\
V & U & $0/55$ & $258/4{,}006$ & & $0$\\
L$_1$ & U & $0/51$ & $26/2{,}930$ & & $0$\\
U & U & $0/59$ & $94/4{,}592$ & & $0$\\
N$_2^B$ & U & $20/202$ & $4{,}856/47{,}449$ & & $22$; $445$\\
N$_2^{N(A)}$ & U & $12/70$ & $2{,}239/13{,}368$ & & $12$; $222$\\
XC & U & $2/46$ & $781/4{,}610$ & & $0$; $28$\\
W & U & $30/120$ & $1{,}858/12{,}712$ & & $48$; $2{,}229$\\
\bottomrule
\end{tabular}}
\end{center}
In the last column, ``$a$; $b$'' gives the false instances at orders $5$ and $7$. The row
N$_1\cup$NA under T$^*$ is the computational shadow of Theorem~\ref{thm:red}: at odd orders, every
instance of N moves into N, into NA, or into a base case.

\section{The obstruction}\label{sec:obs}

\subsection{The single-vertex class}
Under the unconditional calculus, the stuck instances of N$_1$ are \emph{exactly} the all-tight
instances with no other move. Every stuck instance has no target outside $N[y]\cup\{x\}$, and
every instance with such a target moves by~(2) of the proof of Theorem~\ref{thm:red}. There are
$4$ stuck at order $5$, $214$ at order $7$ and $37{,}652$ at order $9$
(\texttt{n1stuck.py}). The smallest are two of order~$5$:
\begin{itemize}
\item \texttt{DEk}: edges $03,04,13,24,34$, $g=0$, $(y,x)=(3,2)$ or $(4,1)$;
\item \texttt{DTw}: edges $02,03,04,14,23,24$, $g=3$, $(y,x)=(0,1)$ or $(2,1)$.
\end{itemize}
In both, $d(g)=2$ and the only target other than $x$ or $y$ is tight. Adding the transfer:
\begin{itemize}
\item with TR the stuck instances are those with $y\in N(g)$ ($154$ at order $7$);
\item with TR$^*$ all of them have $N(g)=\{x,y\}$ ($14$ at order $7$, $1{,}521$ at order $9$);
\item with NA added to the class, none remain.
\end{itemize}
So the all-tight configurations are the obstruction, and a transfer statement through a tight
target is exactly what removes it. The stuck children are good, so this is not a defect of N; it
is the transfer side of~\cite{L1I}.

\subsection{The residual NA}
NA lives in graphs without an even vertex. There the only moves are the leaf graph (d), the split
(g) at degree three, and leaf paths. The leaf graph at the witness $w$ turns $(G,y,x)$ into
$(G-yw+w\ell,\ \ell,\ \{x,y\})$:
\begin{itemize}
\item $y$ is now the even vertex, and $x\in N(y)$;
\item the tail $\ell$ is a leaf at $w$, and $N(w)\cap\{x,y\}=\emptyset$.
\end{itemize}
This is an N$_2$ instance with $g\in A$ and $A\cap N(g)\ne\emptyset$, outside N$_1$.

The smallest NA instances with no move into N$_1\cup$NA are of order $6$: \texttt{EF\char`~w}, which
is $K_6$ minus a triangle $\{0,1,2\}$, with $y=3$ and $x\in\{0,1,2\}$. Adding the set-form instances
with $A\subseteq N[g]$ (class C$_1$) closes the even orders $6$ and $8$. It reopens the odd
orders, with $164$ stuck at order $7$ under TR$^*$. Every one of them has a good child, but none
inside C$_1$. In $89$ of them $g\in A$, and in $102$ every witness is adjacent to $g$.

\subsection{Set forms, far ends, and the triangle at a degree-two even vertex}
\begin{itemize}
\item \emph{Set forms.} The set form N$_2$ never closes: $2$, $249$ and $112{,}226$ stuck at
orders $5$, $7$, $9$. With TR$^*$, $72$ remain at order $7$, each with a good child: $g\in A$ in
$54$ of them, and every witness is adjacent to $g$ in $58$. The union U with the $L_1^*$-admissible instances inherits the blocked
family of~\cite{L1I}: $26$ stuck in L$_1$ alone at order $7$.
\item \emph{Subcubic graphs.} All moves preserve maximum degree, and in a subcubic graph $d(g)=2$.
Even so, N$_2$ does not close: $8$ stuck at order $7$, $182$ at order $9$. All $8$ at order $7$
have $g$ on a triangle.
\item \emph{Far-end classes.} Every far-end class we tried contains false instances. The smallest
is in \texttt{DEk} (edges $03,04,13,24,34$, $g=0$): $c=1$, $A=\{2\}$, $B=\{4\}$. The witness is
$3$, and $N(3)\cap A=\emptyset$.
  \begin{itemize}
  \item $Q_1$ starts $1\,3$. It cannot end at $3$: the remaining edges $03,04,24,34$ would have to
  form one path from $2$ to $4$, with four edges on four vertices.
  \item It cannot end at $4$, as $4\in B$, and it cannot reach $2\in A$.
  \end{itemize}
  The witness condition does not control where $Q_c$ may stop.
\item \emph{The triangle at a degree-two even vertex.} This is the recurring shape. If $d(g)=2$ and
$N(g)=\{u,v\}$ with $u\sim v$, then $g$ subdivides one of two parallel edges $uv$. Suppressing $g$
leaves a multigraph, where Lov\'asz's theorem fails: a triangle with one doubled edge has two odd
vertices and four edges on three vertices. Move~(f) only removes $g$ from $A$.
  \begin{itemize}
  \item The natural reduction deletes $g$ and $uv$, leaving $K=G-g-uv$ with no even vertex.
  Re-inserting the triangle needs $v\notin Q^K(u)$ and $u\notin Q^K(v)$, or disjoint
  $Q^K(u),Q^K(v)$, besides $Q^K(c)\ne Q^K(u),Q^K(v)$.
  \item That is a system of three or four atoms, outside single-tail instances.
  \end{itemize}
  The $7$ instances of N of order~$9$ that the extended calculus leaves uncertified, all $8$ stuck
  subcubic instances at order $7$, and the smallest false far-end instances all have this shape.
\end{itemize}

\subsection*{What persists under every refinement}
\begin{enumerate}
\item The \emph{all-tight configurations}. Only a transfer through a tight target ($\TRN$, or TR
of~\cite{L1I}) removes them; no instance-level move into a single-tail class does.
\item \emph{NA}. Its leaf-graph children carry a forbidden pair $\{x,g\}$ with $x\in N(g)$, and no
single-tail class containing them closed.
\item \emph{The degree-two even vertex on a triangle}, where the natural reduction needs a
multi-atom system.
\end{enumerate}

\section{Formalisation}\label{sec:lean}

\texttt{lean/GallaiLean/ConjN.lean} and \texttt{ConjNFar.lean} (Lean 4, Mathlib) have no
\texttt{sorry}. Every theorem below depends only on \texttt{propext}, \texttt{Classical.choice} and
\texttt{Quot.sound}, and the full library builds.

\begin{center}\small
\resizebox{\textwidth}{!}{\begin{tabular}{ll}
\toprule
statement & Lean\\
\midrule
Conjecture~\ref{conj:n}, $\TRN$, NA, degree-two case & \texttt{ConjNOn}, \texttt{TightTransferNOn}, \texttt{NeighbourAvoidOn}, \texttt{ConjNDegTwoOn}\\
Lemma~\ref{lem:split} & \texttt{avoid\_of\_splitOff}\\
Lemma~\ref{lem:deg2} & \texttt{conjNDegTwo\_of\_neighbourAvoid}\\
Theorem~\ref{thm:red} & \texttt{conjN\_of\_tight\_degTwo} (induction \texttt{conjN\_strong}), \texttt{conjN\_of\_tight\_neighbourAvoid}\\
Corollary~\ref{cor:TR} & \texttt{tightTransferNOn\_of}, \texttt{conjN\_of\_TR\_neighbourAvoid}; \texttt{TightTransferAdj}\\
Lemma~\ref{lem:protected}, ends outside $N_K(u)$ & \texttt{fan\_single\_protected\_path} (with \texttt{PathDecomp.erase\_path}, \texttt{PathDecomp.add\_path})\\
Theorem~\ref{thm:far}, move (x) & \texttt{avoidFar\_target}; extended instances \texttt{AvoidFarDecomp}\\
\bottomrule
\end{tabular}}
\end{center}

\begin{itemize}
\item \emph{Library results used.} \texttt{lovasz\_avoid\_nonadj} ($L_1$),
\texttt{avoid\_free\_target} (Case~I, via \texttt{csys\_caseI}), \texttt{splitOff\_facts},
\texttt{exists\_subst\_path}, \texttt{exists\_decomp\_of\_edgesM'}, and the definitions
\texttt{AvoidDecomp}, \texttt{TightTransfer}.
\item \emph{Not formalised.} Lemma~\ref{lem:leafpath}; the degree-one case of
Lemma~\ref{lem:protected}; the moves (b$_B$), (C), (c$_B$) of Theorem~\ref{thm:far}. These are
used only in the computations of Section~\ref{sec:comp}.
\end{itemize}

\section{What remains}\label{sec:open}
\begin{enumerate}
\item $\TRN$, the tight transfer for one forbidden vertex. For $y\notin N(g)$ this is TR
of~\cite{L1I} with $|A|=1$; for $y\in N(g)$, its counterpart. Partial cases are listed
in~\cite[\S4.1]{L1I}.
\item NA, neighbour avoidance in graphs without an even vertex, when $x,y$ have a common
neighbour. By Lemma~\ref{lem:deg2} it is the degree-two case of~N.
\item An inductive class that handles the triangle at a degree-two even vertex. This most likely
needs systems of two or three atoms in the sense of~\cite{CLT}, not single-tail instances.
\item The strengthening $N^+$ for $x\sim g$, which held in every case to order $9$. Its far-end
constraint is what move~(x) needs. It would be a natural hypothesis for a far-end class that
handles the core.
\end{enumerate}

\begin{thebibliography}{9}
\bibitem{CLT} Constrained Lov\'asz decompositions, \texttt{manuscripts/clt}.
\bibitem{DL} The disjunctive avoidance conjecture for one pair, \texttt{manuscripts/dl-one}.
\bibitem{Fan05} G.~Fan, Path decompositions and Gallai's conjecture, \emph{J. Combin. Theory Ser.
B} 93 (2005) 117--125.
\bibitem{L1I} Lov\'asz's induction for $L_1^*$, \texttt{manuscripts/l1-induction}.
\bibitem{L1S} The avoidance conjecture $L_1^*$, \texttt{manuscripts/l1star}.
\bibitem{L68} L.~Lov\'asz, On covering of graphs, in: \emph{Theory of Graphs} (Proc. Colloq.,
Tihany, 1966), Academic Press, 1968, 231--236.
\bibitem{MS} Multi-pair separation, \texttt{manuscripts/multi-sep}.
\bibitem{RD} Route descent, \texttt{manuscripts/route-descent}.
\bibitem{STL} Star-terminal Lov\'asz decompositions, \texttt{manuscripts/stl}.
\end{thebibliography}
\end{document}
